\section{Degree-eight vertices without computation}\label{sec:degree-eight}

\begin{lemma}\label{lem:independent-triple}
Let $G$ satisfy the hypotheses of Theorem~\ref{thm:main}.  Every
degree-eight vertex $v$ satisfies $\alpha(G[N(v)])=3$.
\end{lemma}

\begin{proof}
Dirac's inequality (3.1) gives the upper bound three.  Put $J=G[N(v)]$.
By Proposition~\ref{prop:critical-base}, $J$ has no $K_4$ subgraph and
$|V(G)|\geq n_8\geq25$, so $G-N[v]$ is nonempty.  If $\alpha(J)\leq2$,
Rolek, Song and Thomas~\cite[Lemma~2.1]{RolekSongThomas2022} show that $J$
contains a spanning square of an eight-cycle.  Label its vertices
cyclically $0,\ldots,7$.  An exterior component has at least seven
neighbours in $N(v)$, by seven-connectivity.  Contract that component to
$c$ and delete the other exterior vertices.  By cyclic symmetry and edge
deletion, take $c$ adjacent to every cycle vertex except $0$.  The sets
\[
       \{0,7,2\},\quad\{3\},\quad\{4\},\quad\{1,v\},
       \quad\{6\},\quad\{5\},\quad\{c\}
\]
are connected and pairwise adjacent except for $\{3\}:\{6\}$.  They
form a forbidden $K_7^{-}$ model.  Hence $\alpha(J)=3$.
\end{proof}

\begin{lemma}\label{lem:good-deletion}
Let $J$ be an eight-vertex graph with minimum degree at least four, no
$K_4$ subgraph, and independence number three.  Call $r$ good if $J-r$
has a $K_5^{-}$ minor.  Every vertex of $J$ has a good neighbour.
\end{lemma}

\begin{proof}
The graph $J$ is three-connected.  Indeed, a component has order at least
five, and a component behind a cutvertex has order at least four, excluding
both kinds of separation on eight vertices.  A two-cut would leave two
three-vertex components, each a triangle complete to the cut.  An edge in
the cut gives a $K_4$; an independent cut gives independence number two.
Both are impossible.

Suppose first that $J$ is four-connected.  Then $J-r$ is three-connected
for every $r$.  Wood and Woodall~\cite[Lemma~4.2.1]{WoodWoodall2009}
classify three-connected $K_5^{-}$-minor-free graphs as wheels, the
triangular prism and $K_{3,3}$.  If some $r$ is not good, order seven
forces $J-r$ to be a wheel with six rim vertices.  Every rim vertex must
see $r$ to have degree at least four, and the hub cannot see $r$ because
$J$ has no $K_4$.  Thus $J=\overline{K_2}\vee C_6$.  With apices $p,q$
and rim $v_0,\ldots,v_5$, deletion of $v_2$ leaves the $K_5^{-}$ model
\[
       \{p,v_0\},\quad\{q,v_1\},\quad\{v_3\},
       \quad\{v_4\},\quad\{v_5\},
\]
missing only $\{v_3\}:\{v_5\}$.  Every rim vertex is therefore good,
and the rim vertices meet every neighbourhood.  If no bad $r$ exists,
the conclusion is immediate.

Now let $S$ be a three-cut.  Its two components have orders two and three,
since each has order at least two.  Write the first as $A=\{a,a'\}$ and
the second as $B$.  The edge $aa'$ is present and both ends are complete
to $S$.  Consequently $S$ is independent.  Every vertex of $B$ has a
neighbour in $B$, so $J[B]$ is a path or triangle.

If $B=b_1b_2b_3$ is a path, its ends are complete to $S$ and $b_2$ has at
least two neighbours in $S$.  The following models establish that all
vertices in $A\cup B$ are good; symmetry handles $a'$ and $b_3$.
In the table, $b_1s_1$ abbreviates the branch set $\{b_1,s_1\}$, and
similarly for the other entries.
\[
\begin{array}{c|c|c}
\text{deleted vertex}&\text{five branch sets}&\text{only missing pair}\\\hline
a&b_1s_1,b_3s_2,b_2,a',s_3&b_2:a'\\
b_1&as_1,b_3s_2,b_2,a',s_3&b_2:a'\\
b_2&as_1,b_1s_2,b_3,a',s_3&b_3:a'
\end{array}
\]
Here $S=\{s_1,s_2,s_3\}$ is labelled with $b_2s_3$ present in the first
row, and with $b_2s_1,b_2s_3$ present in the second.

If $B$ is a triangle, each of its vertices sees at least two vertices of
$S$, and no vertex of $S$ sees all of $B$.  Its missing incidences with
$S$ therefore form a perfect matching; label these $b_is_i$.
After deleting $a$ use
\[
       s_1b_2,\quad s_2a',\quad b_1,\quad s_3,\quad b_3,
\]
which misses only $s_3:b_3$; after deleting $b_1$ use
\[
       s_1b_2b_3,\quad s_2,\quad s_3,\quad a,\quad a',
\]
which misses only $s_2:s_3$.  Symmetry again makes all vertices in
$A\cup B$ good.  In either case, the vertices of $A$ see one another,
every vertex of $B$ has a neighbour in $B$, and every vertex of $S$ sees
$A$.  Thus every vertex has a good neighbour.
\end{proof}

\begin{proposition}\label{prop:critical-three-edge}
Let $G$ satisfy the hypotheses of Theorem~\ref{thm:main}.  Every
degree-eight vertex is incident with an edge $e$ satisfying $c(e)\leq3$,
and $n_8\geq26+\tau$.  Both conclusions have computation-free proofs.
\end{proposition}

\begin{proof}
Fix a degree-eight vertex $v$ and put $J=G[N(v)]$.  By
Proposition~\ref{prop:critical-base} and Lemma~\ref{lem:independent-triple},
$J$ is $K_4$-free and has independence number three.  If every incident
edge has codegree at least four, then $\delta(J)\geq4$.  An exterior
component exists because $|V(G)|\geq25$, and seven-connectivity gives at
least seven boundary neighbours.  Contract it to $c$ and delete the other
exterior components.  Let $x$ be the vertex of $J$ missed by $c$, or any
vertex if none is missed.  By Lemma~\ref{lem:good-deletion}, $x$ has a
neighbour $r$ such that $J-r$ has a $K_5^{-}$ model $M_1,\ldots,M_5$.
Then $cr$ is an edge, and
\[
                  \{v\},\quad\{c,r\},\quad M_1,\ldots,M_5
\]
form a $K_7^{-}$ model: the first two sets are adjacent through $vr$ and
are adjacent to every $M_i$.  The sole possible missing contact from $c$
is a singleton $M_i=\{x\}$, supplied instead by $rx$.  This contradiction
proves that some edge $e=vx$ has $c(e)=d_J(x)\leq3$.

The quotient $G/e$ is six-connected, has no $K_7^{-}$ minor and has order
at least twenty-four.  Lemma~\ref{lem:jakobsen-defect} gives $D(G/e)\geq25$,
while \eqref{eq:contraction} gives $D(G)\geq D(G/e)+1\geq26$.
The identity $D(G)=n_8-\tau$ proves the bound.
\end{proof}
